\documentclass[12pt, titlepage]{article}

\usepackage{booktabs}
\usepackage{tabularx}
\usepackage{amsmath}
\usepackage{hyperref}
\hypersetup{
    colorlinks,
    citecolor=black,
    filecolor=black,
    linkcolor=red,
    urlcolor=blue
}
\usepackage[round]{natbib}

\input{../Comments.text}
\input{../Common.text}

\begin{document}

\title{Verification and Validation Report: \progname}
\author{\authname}
\date{April 15, 2026}

\maketitle

\pagenumbering{roman}

\section{Revision History}

\begin{tabularx}{\textwidth}{p{3cm}p{2cm}X}
\toprule {\bf Date} & {\bf Version} & {\bf Notes}\\
\midrule
Apr 15, 2026 & 1.0 & Initial VnV Report\\
Apr 17, 2026 & 1.1 & Updated after code review: fixed DetectorHit interpolation and Control call order\\
\bottomrule
\end{tabularx}

~\newpage

\section{Symbols, Abbreviations and Acronyms}

\renewcommand{\arraystretch}{1.2}
\begin{tabular}{l l}
  \toprule
  \textbf{Symbol} & \textbf{Description}\\
  \midrule
  T    & Test\\
  FR   & Functional Requirement\\
  INV  & Input Validation\\
  ODE  & Ordinary Differential Equation\\
  RK45 & Runge--Kutta 4(5) (Dormand--Prince)\\
  \bottomrule
\end{tabular}

\newpage

\tableofcontents

\listoftables

\newpage

\pagenumbering{arabic}

This report documents the VnV results for Trajecto, a 2-D
charged-particle trajectory simulator generated by the Drasil framework.
Tests were run on the Python implementation at
\texttt{code/drasil-example/trajecto/src/python/}.
All test input files and scripts are in \texttt{CAS741/test/}.

% -----------------------------------------------------------------------
\section{Functional Requirements Evaluation}
% -----------------------------------------------------------------------

Six functional tests (T-FR-1 to T-FR-6) were run.
Table~\ref{tab:fr-results} shows the results.

\begin{table}[h]
\caption{Functional Requirement Test Results}
\label{tab:fr-results}
\centering
\begin{tabular}{lllp{7cm}}
\toprule
\textbf{Test} & \textbf{Script} & \textbf{Result} & \textbf{Notes} \\
\midrule
T-FR-1 & test\_fr1.py & PASS &
  All 17 input fields were echoed correctly.
  Soft warnings for $m$, $q$, $t_\text{final}$ being out of the
  physical range are expected because the test uses simple
  round numbers to focus on the maths. \\
T-FR-2 & test\_fr2.py & PASS &
  No-field case. Analytic solution: $x(t)=2t$, $y(t)=t$.
  Max position error: $1.33\times10^{-15}$\,m (limit: $10^{-4}$\,m).
  Vertical detector hit at $t=0.5$\,s, $x=1.0$\,m, $y=0.5$\,m. \\
T-FR-3 & test\_fr3.py & PASS &
  Magnetic field only. Speed should not change.
  Max relative speed-squared drift: $1.11\times10^{-15}$ (limit: 1\%).
  Detector returns $t_\text{hit}=-1$ correctly (particle never hits). \\
T-FR-4 & test\_fr4.py & PASS &
  Constant electric field $E_x=1$\,N/C. Analytic: $x(t)=\tfrac{1}{2}t^2$.
  Max position error: $7.15\times10^{-14}$\,m (limit: $10^{-3}$\,m).
  Comparison is limited to $x\le2$\,m where the field is active;
  beyond that point the particle leaves the grid and the analytic
  formula no longer applies.
  Hit: $t\approx2.000$\,s, $x=2.000$\,m, $y=0$\,m (interpolated to exact detector position). \\
T-FR-5 & test\_fr5.py & PASS &
  Particle moves through two regions: first field-free, then magnetic.
  Hit confirmed at $(0.5,\,0.0)$\,m.
  Speed drift in Region~2: $1.51\times10^{-8}$ (limit: 1\%),
  confirming the field switched correctly. \\
T-FR-6 & test\_fr6.py & PASS &
  Horizontal detector hit at $t=1.0$\,s, $x=1.0$\,m, $y=1.0$\,m. \\
\bottomrule
\end{tabular}
\end{table}

\subsection{Input Validation Tests (Section~4.1.7)}

Table~\ref{tab:inv-results} shows the results for T-INV-1 to T-INV-12.

\begin{table}[h]
\caption{Input Validation Test Results}
\label{tab:inv-results}
\centering
\begin{tabular}{lllp{7cm}}
\toprule
\textbf{Test} & \textbf{Input} & \textbf{Result} & \textbf{Notes} \\
\midrule
T-INV-1 & $m=0$ & PASS &
  \texttt{InputError} is raised and the program exits with code~1. \\
T-INV-2 & $q=0$ & KNOWN GAP &
  No $q>0$ check exists in the generated code.
  The program runs with $\kappa=0$ (no Lorentz force, straight-line
  motion) and exits with code~0.
  The VnV Plan expected an error here.
  This is a gap in the Drasil-generated input constraint module. \\
T-INV-3 & $t_\text{final}=-1$ & PASS &
  \texttt{InputError} raised with a message about \texttt{t\_final}.
  Exits with code~1. \\
T-INV-4 & $d_\text{orient}=2$ & KNOWN GAP &
  The code only prints a soft warning and exits with code~0.
  The VnV Plan expected a hard error and a non-zero exit code.
  This is a gap in the Drasil-generated input constraint module. \\
T-INV-5 & $d_\text{len}=0$ & PASS &
  \texttt{InputError} raised with a message about \texttt{d\_len}.
  Exits with code~1. \\
T-INV-6--12 & various & NOT APPLICABLE &
  These tests require region shapes that cannot be expressed in the
  grid-based input format.
  The code always produces uniform, axis-aligned, gap-free tiles,
  so invalid geometry configurations are impossible to create. \\
\bottomrule
\end{tabular}
\end{table}

\subsection{Coverage Summary (Section~4.1.8)}

Table~\ref{tab:cov-constraint} shows which input constraints were tested.
Table~\ref{tab:cov-oracle} shows the pseudo-oracle and output-contract
coverage.

\begin{table}[h]
\caption{Input-Constraint Coverage}
\label{tab:cov-constraint}
\centering
\begin{tabular}{llll}
\toprule
\textbf{Constraint} & \textbf{Expected Behaviour} & \textbf{Test} & \textbf{Status} \\
\midrule
$m>0$                      & Error + exit 1       & T-INV-1  & PASS \\
$q>0$                      & Error + exit 1       & T-INV-2  & KNOWN GAP \\
$t_\text{final}>0$         & Error + exit 1       & T-INV-3  & PASS \\
$d_\text{orient}\in\{0,1\}$& Error + exit 1       & T-INV-4  & KNOWN GAP \\
$d_\text{len}>0$           & Error + exit 1       & T-INV-5  & PASS \\
Detector location          & N/A (grid input)     & T-INV-6  & N/A \\
Region geometry            & N/A (grid input)     & T-INV-7,8& N/A \\
Equal region sizes         & N/A (grid enforces)  & T-INV-9  & N/A \\
Adjacency                  & N/A (grid tiles)     & T-INV-10 & N/A \\
No overlap                 & N/A (grid tiles)     & T-INV-11 & N/A \\
Union is rectangle         & N/A (grid union)     & T-INV-12 & N/A \\
\bottomrule
\end{tabular}
\end{table}

\begin{table}[h]
\caption{Pseudo-Oracle and Output-Contract Coverage}
\label{tab:cov-oracle}
\centering
\begin{tabular}{llll}
\toprule
\textbf{Category} & \textbf{Reference} & \textbf{Tests} & \textbf{Status} \\
\midrule
No-field analytic     & $x=2t$,\ $y=t$             & T-FR-2          & PASS \\
$E$-only analytic     & $x=\tfrac{1}{2}t^2$         & T-FR-4          & PASS \\
$B$-only speed check  & $|\Delta v^2|/v_0^2\le1\%$  & T-FR-3, T-FR-5  & PASS \\
Multi-region switch   & Speed stays constant in R2  & T-FR-5          & PASS \\
No-hit sentinel       & $t_\text{hit}=-1$           & T-FR-3          & PASS \\
Hit values reported   & $t_\text{hit}\ge0$, $(x,y)$ & T-FR-2,4,5,6    & PASS \\
\bottomrule
\end{tabular}
\end{table}

% -----------------------------------------------------------------------
\section{Nonfunctional Requirements Evaluation}
% -----------------------------------------------------------------------

\subsection{Accuracy}

Accuracy was checked using analytic solutions as pseudo-oracles.

For the no-field case (T-FR-2), the analytic trajectory is $x(t)=2t$,
$y(t)=t$.
The maximum position error from the simulation was $1.33\times10^{-15}$\,m,
which is effectively machine precision.

For the constant electric-field case (T-FR-4), the analytic trajectory is
$x(t)=\tfrac{1}{2}t^2$.
The maximum error was $7.15\times10^{-14}$\,m.

For magnetic-field cases (T-FR-3, T-FR-5), speed must stay constant because
the magnetic force does no work.
The maximum relative drift in $v^2$ was $1.51\times10^{-8}$, which is
far below the 1\% limit set in the VnV Plan.

Overall, the simulator meets the accuracy requirement.

\subsection{Usability}

Trajecto is a command-line tool.
The user runs it with one command: \texttt{python3 Control.py input.txt}.
The input file follows a fixed format with one label line before each value.

The program prints a clear \texttt{InputError} message when required fields
are invalid (e.g.\ $m\le0$, $t_\text{final}\le0$, $d_\text{len}\le0$).
For values outside the physical range, a \texttt{Warning} is printed so
the user knows the input is unusual but the run still completes.
Output is written to \texttt{output.txt} in a readable key-value format.

The interface is simple and consistent.
No GUI testing was performed, as Trajecto has no graphical interface.

\subsection{Maintainability}

Trajecto is generated automatically by the Drasil framework.
The source code is not meant to be edited directly.
Changes to the system are made by updating the Drasil specification and
regenerating the code.

The generated Python code is clean and follows a standard module layout:
\texttt{Control}, \texttt{InputParameters}, \texttt{Calculations},
\texttt{DetectorHit}, \texttt{OutputFormat}, and \texttt{Constants}.
Each module has a single clear responsibility.

Because the code is generated, manual maintenance effort is minimal.
The main maintainability concern is keeping the Drasil specification up to
date, which is outside the scope of this VnV activity.

\subsection{Portability}

The Python implementation was tested on macOS using Python~3.
It depends only on the standard library and \texttt{scipy}, which are
available on all major platforms (macOS, Linux, Windows).

The test scripts in \texttt{CAS741/test/} use only standard Python modules
(\texttt{subprocess}, \texttt{os}, \texttt{ast}, \texttt{math}) and require
no additional installation.
Any machine with Python~3 and \texttt{scipy} installed can run the simulator
and all tests without modification.

% -----------------------------------------------------------------------
\section{Comparison to Existing Implementation}
% -----------------------------------------------------------------------

Not applicable. Trajecto is a new artefact with no prior reference
implementation to compare against.

% -----------------------------------------------------------------------
\section{Unit Testing}
% -----------------------------------------------------------------------

% -----------------------------------------------------------------------
\section{Changes Due to Testing}
% -----------------------------------------------------------------------

Two bugs were found and fixed in the Drasil specification (Haskell source),
and the code was regenerated:
\begin{itemize}
  \item \textbf{Call order in Control.py}: \texttt{genAllInputCalls} in
        \texttt{drasil-code} generated \texttt{derived\_values()} before
        \texttt{input\_constraints()}, causing a \texttt{ZeroDivisionError}
        when $m=0$ instead of the expected \texttt{InputError}.
        Fixed by swapping the order of \texttt{genConstraintCall} and
        \texttt{genDerivedCall} in \texttt{FunctionCalls.hs}.
  \item \textbf{Interpolation in DetectorHit.py}: The hit position and time were
        recorded at the ODE step after the crossing, not at the crossing itself.
        Fixed in \texttt{ModuleDefs.hs} by adding linear interpolation: the
        fractional step \texttt{frac} is computed, and the hit coordinate and
        time use \texttt{(i - 1 + frac) * t\_step} and the interpolated
        position.
        This also required fixing an operator-precedence error in the Haskell DSL
        expression (\texttt{\$+} and \texttt{\$*} share default fixity
        \texttt{infixl 9}), resolved by adding explicit parentheses.
\end{itemize}

Two remaining implementation gaps were found and documented:
\begin{itemize}
  \item \textbf{T-INV-2 ($q=0$)}: The generated \texttt{InputParameters.py}
        has no hard check for $q>0$.
        To fix this, a constraint \texttt{raise Exception("InputError")} for
        $q \le 0$ should be added to \texttt{input\_constraints()} in the
        Drasil specification, and the code should be regenerated.
  \item \textbf{T-INV-4 ($d_\text{orient} \notin \{0,1\}$)}: The generated
        code only prints a soft warning.
        To fix this, the check for \texttt{d\_orient} should be changed from
        a warning to a hard error in the Drasil specification, and the code
        should be regenerated.
\end{itemize}
Seven tests (T-INV-6--12) were marked not applicable.
The VnV Plan assumed a corner-based region input format, but the generated
code uses a uniform grid.
The test scripts were updated to match the actual code behaviour.

% -----------------------------------------------------------------------
\section{Automated Testing}
% -----------------------------------------------------------------------

All tests are Python scripts in \texttt{CAS741/test/}.
Each script calls \texttt{Control.py} through \texttt{subprocess},
reads \texttt{output.txt}, and prints a PASS or FAIL verdict.
To run all tests:
\begin{verbatim}
cd CAS741/test
python3 test_fr1.py
python3 test_fr2.py
python3 test_fr3.py
python3 test_fr4.py
python3 test_fr5.py
python3 test_fr6.py
python3 test_inv.py
python3 test_coverage.py
\end{verbatim}

% -----------------------------------------------------------------------
\section{Trace to Requirements}
% -----------------------------------------------------------------------

The requirements are: R1 (input), R2 (trajectory), R3 (detector), R4
(multi-region field), R5 (output), NFR1 (accuracy), NFR2 (usability),
NFR3 (maintainability), NFR4 (portability).

\begin{table}[ht]
\caption{Traceability Between Test Cases and Requirements}
\label{tab:trace-req}
\centering
\begin{tabular}{|c|c|c|c|c|c|c|c|c|c|}
\hline
         & R1 & R2 & R3 & R4 & R5 & NFR1 & NFR2 & NFR3 & NFR4 \\
\hline
T-FR-1   & X  &    &    &    & X  &      &      &      &      \\
\hline
T-FR-2   &    & X  & X  &    & X  & X    &      &      &      \\
\hline
T-FR-3   &    & X  & X  &    & X  & X    &      &      &      \\
\hline
T-FR-4   &    & X  & X  &    & X  & X    &      &      &      \\
\hline
T-FR-5   &    & X  & X  & X  & X  & X    &      &      &      \\
\hline
T-FR-6   &    & X  & X  &    & X  &      &      &      &      \\
\hline
T-INV-1  & X  &    &    &    &    &      &      &      &      \\
\hline
T-INV-2  & X  &    &    &    &    &      &      &      &      \\
\hline
T-INV-3  & X  &    &    &    &    &      &      &      &      \\
\hline
T-INV-4  & X  &    &    &    &    &      &      &      &      \\
\hline
T-INV-5  & X  &    &    &    &    &      &      &      &      \\
\hline
Test Accuracy      &    &    &    &    &    & X    &      &      &      \\
\hline
Test Usability     &    &    &    &    &    &      & X    &      &      \\
\hline
Test Maintainability &  &    &    &    &    &      &      & X    &      \\
\hline
Test Portability   &    &    &    &    &    &      &      &      & X    \\
\hline
\end{tabular}
\end{table}

% -----------------------------------------------------------------------
\section{Trace to Modules}
% -----------------------------------------------------------------------

The modules are: M1~=~Control, M2~=~InputParameters, M3~=~Calculations,
M4~=~DetectorHit, M5~=~OutputFormat, M6~=~Constants.

\begin{table}[ht]
\caption{Traceability Between Test Cases and Modules}
\label{tab:trace-mod}
\centering
\begin{tabular}{|c|c|c|c|c|c|c|}
\hline
         & M1 & M2 & M3 & M4 & M5 & M6 \\
\hline
T-FR-1   & X  & X  &    &    & X  & X  \\
\hline
T-FR-2   & X  &    & X  & X  & X  & X  \\
\hline
T-FR-3   & X  &    & X  & X  & X  & X  \\
\hline
T-FR-4   & X  &    & X  & X  & X  & X  \\
\hline
T-FR-5   & X  &    & X  & X  & X  & X  \\
\hline
T-FR-6   & X  &    & X  & X  & X  & X  \\
\hline
T-INV-1  & X  & X  &    &    &    &    \\
\hline
T-INV-2  & X  & X  &    &    &    &    \\
\hline
T-INV-3  & X  & X  &    &    &    &    \\
\hline
T-INV-4  & X  & X  &    &    &    &    \\
\hline
T-INV-5  & X  & X  &    &    &    &    \\
\hline
\end{tabular}
\end{table}

% -----------------------------------------------------------------------
\section{Code Coverage Metrics}
% -----------------------------------------------------------------------

Formal coverage metrics were not collected.
Trajecto is generated code; coverage of the generator itself is handled by
the Drasil project's own tests.
The functional test suite runs code in all six modules (Control,
InputParameters, Calculations, DetectorHit, OutputFormat, Constants) and
covers every scenario reachable through the input format: field-free,
$E$-only, $B$-only, mixed, hit, no-hit, and hard constraint violations.

% -----------------------------------------------------------------------
\section{Extras}
% -----------------------------------------------------------------------

No extras were planned for this deliverable.

\bibliographystyle{plainnat}
\bibliography{../../refs/References}

\newpage{}
\section*{Appendix --- Reflection}

The information in this section will be used to evaluate the team members on the
graduate attribute of Reflection.

\input{../Reflection.text}

\begin{enumerate}
  \item What went well while writing this deliverable?
  \item What pain points did you experience during this deliverable, and how
        did you resolve them?
  \item Which parts of this document stemmed from speaking to your client(s) or
        a proxy (e.g. your peers)? Which ones were not, and why?
  \item In what ways was the Verification and Validation (VnV) Plan different
        from the activities that were actually conducted for VnV?  If there were
        differences, what changes required the modification in the plan?  Why did
        these changes occur?  Would you be able to anticipate these changes in
        future projects?  If there weren't any differences, how was your team
        able to clearly predict a feasible amount of effort and the right tasks
        needed to build the evidence that demonstrates the required quality?
        (It is expected that most teams will have had to deviate from their
        original VnV Plan.)
\end{enumerate}

\end{document}%
